\section{The arithmetic positivity target}\label{sec:arithmetic}

This section fixes the target independently of any field-theory model.
The formulas also explain why positivity of a single covariance kernel
would be insufficient.

\subsection{The compactly supported Weil form}

Let $I_L=(-L/2,L/2)$, with $L>0$, and let
$f\in C_c^\infty(I_L;\C)$. Write $F$ for its zero extension to $\R$ and use
\begin{equation}
 \widehat F(\tau)=\int_\R e^{-\ii\tau x}F(x)\dd x,
 \qquad \ip f g=\int_{I_L}\overline{f(x)}g(x)\dd x.
\end{equation}
For $a>0$, define the compressed shift $T_af(x)=F(x-a)$ on $I_L$.
Let $\psi=\Gamma'/\Gamma$ be the digamma function, and let
$\Lambda(n)=\log p$ when $n=p^k$, $p$ prime and $k\geq1$, and zero otherwise.
Finally put
\begin{equation}
 C(f)=\int_{I_L}f(x)\cosh(x/2)\dd x,
 \qquad S(f)=\int_{I_L}f(x)\sinh(x/2)\dd x.
\end{equation}
The normalization of the program is
\begin{align}\label{eq:weil}
 Q_{0,L}[f]={}&\frac1{2\pi}\int_\R
 \bigl[\Ree\psi(\tfrac14+\tfrac{\ii\tau}2)-\log\pi\bigr]
 |\widehat F(\tau)|^2\dd\tau
 +2|C(f)|^2-2|S(f)|^2\notag\\
 &-\sum_{\substack{n\geq2\\\log n<L}}
 \frac{\Lambda(n)}{\sqrt n}\ip f{(T_{\log n}+T_{\log n}^*)f}.
\end{align}
The prime-power sum is finite on a fixed interval. The pole term is
indefinite, and the archimedean multiplier includes its local constant.
Neither may be dropped in a positivity argument.

For comparison with the usual explicit formula, set
$h=F*\widetilde F$, where $\widetilde F(x)=\overline{F(-x)}$.
The two prime sums in the Weil functional give the last line of
\eqref{eq:weil}; its $e^{x/2}+e^{-x/2}$ term gives
$2|C|^2-2|S|^2$. Its archimedean distribution has Fourier multiplier
$\Ree\psi(1/4+\ii\tau/2)-\log\pi$. Thus
$Q_{0,L}[f]=W(F*\widetilde F)$ in the convention of
\cite[p.~1]{Suzuki}. In particular the weight is $\Lambda(n)/\sqrt n$.

\begin{theorem}[Weil positivity criterion, used as background]\label{thm:weil}
RH holds if and only if $Q_{0,L}[f]\geq0$ for every $L>0$ and every
$f\in C_c^\infty(I_L;\C)$.
\end{theorem}
The equivalence is the classical criterion, not a result of the present
construction. Every compact support lies in one such interval, so an
unbounded sequence of interval lengths suffices. A finite collection of
lengths or numerical matrices does not.

\subsection{The positive kinetic term and the indispensable local term}

The numbers $a_n=2n+1/2$, $n\geq0$, are the decay rates
$1/2,5/2,9/2,\ldots$ in the exponential expansion
\begin{equation}\label{eq:intro-tower}
 \nGamma(u)=\frac{e^{-u/2}}{1-e^{-2u}}
 =\sum_{n=0}^\infty e^{-a_nu},\qquad u>0.
\end{equation}
They arise by expanding a geometric series, and are not zeta zeros
or prime-power coefficients. In the free defect construction they
also appear as the even angular rates $\ell+1/2$ with $\ell=2n$;
the physical derivation is in \cite[Section~S2]{Supplement}.
Put
\begin{equation}\label{eq:b-symbol}
 b(\tau^2)=\Ree\psi(\tfrac14+\tfrac{\ii\tau}2)-\psi(\tfrac14)
 =2\sum_{n=0}^\infty\frac{\tau^2}{a_n(a_n^2+\tau^2)}\geq0,
 \qquad w_0=\psi(\tfrac14)-\log\pi<0.
\end{equation}
The exact digamma value is
\begin{equation}
 \psi(\tfrac14)=-\gamma-\frac{\pi}{2}-3\log2,\qquad
 w_0=-\gamma-\frac{\pi}{2}-3\log2-\log\pi
       \simeq-5.37218342,
\end{equation}
where $\gamma$ is Euler's constant. Thus $\psi(1/4)$ is not equal
to $\log\pi$. Adding and
subtracting $\psi(1/4)$ in the original multiplier gives the exact identity
\begin{align}\label{eq:archimedean-split}
 \Ree\psi(\tfrac14+\tfrac{\ii\tau}2)-\log\pi
 &=[\Ree\psi(\tfrac14+\tfrac{\ii\tau}2)-\psi(\tfrac14)]
   +[\psi(\tfrac14)-\log\pi]\notag\\
 &=b(\tau^2)+w_0.
\end{align}
Thus $w_0$ is the remainder after separating the nonnegative multiplier
$b(\tau^2)$; it is already present in \eqref{eq:weil}, not an additional
term introduced into the form.

The series in \eqref{eq:b-symbol} follows by subtracting the convergent digamma partial
fractions. Each exponential in \eqref{eq:intro-tower} satisfies
$\int_0^\infty e^{-a u}(1-\cos\tau u)\dd u
=\tau^2/[a(a^2+\tau^2)]$. Plancherel therefore gives
\begin{equation}\label{eq:kinetic}
 \mathcal E_\Gamma[F]
 :=\int_0^\infty\nGamma(u)\norm{F(\,\cdot+u)-F}_2^2\dd u
 =\frac1{2\pi}\int_\R b(\tau^2)|\widehat F(\tau)|^2\dd\tau.
\end{equation}
Using \eqref{eq:archimedean-split} and
$\frac1{2\pi}\int_\R|\widehat F(\tau)|^2\dd\tau=\norm f_2^2$,
the archimedean contribution becomes
\begin{equation}\label{eq:archimedean-energy}
 \frac1{2\pi}\int_\R
 [\Ree\psi(\tfrac14+\tfrac{\ii\tau}2)-\log\pi]
 |\widehat F(\tau)|^2\dd\tau
 =\mathcal E_\Gamma[F]+w_0\norm f_2^2.
\end{equation}
Substituting this identity into \eqref{eq:weil} rewrites the full target as
\begin{align}\label{eq:weil-energy}
 Q_{0,L}[f]={}&\mathcal E_\Gamma[F]+w_0\norm f_2^2
 +2|C(f)|^2-2|S(f)|^2\notag\\
 &-\sum_{\substack{n\geq2\\\log n<L}}
 \frac{\Lambda(n)}{\sqrt n}\ip f{(T_{\log n}+T_{\log n}^*)f}.
\end{align}
Here $w_0\norm f_2^2=w_0\int_{I_L}|f(x)|^2\dd x$ is the fixed
negative local term. The positive difference energy is only one part of
the desired form: its local term, pole contribution and prime-power
translations must all be retained.

The natural form domain on a fixed interval is given by zero extensions
with $\int\log(2+|\tau|)|\widehat F(\tau)|^2\dd\tau<\infty$.
Indeed $b(\tau^2)$ grows logarithmically at infinity, whereas the pole
and finite translation terms are bounded forms on $L^2(I_L)$. This
explains the logarithmic form domain and semiboundedness of the
localized form. For the logical proof target in Theorem~\ref{thm:weil},
the smooth compactly supported class is enough.
